\section{Introduction}
\label{sec:introduction}

Foundation models have revolutionized the way we think about computing. They arose as a paradigm only five years ago \citep{Bommasani2021FoundationModels}, yet in that time have become the basis for new discoveries across a broad set of scientific fields \citep{Jumper2021HighlyAccurate, Lam2023LearningSkillful, Merchant2023ScalingDeep, Trinh2024SolvingOlympiad, Hubert2026Olympiad}. As you might expect, those five years saw a flurry of research into how best to build these models, and one key ingredient emerged: scale. Unfortunately, scale makes experiments incredibly expensive.


The natural solution is to experiment at the small scale, then transfer to the large scales. Despite the early promise of scaling laws \citep{Kaplan2020ScalingLaws}, this solution has not yet materialized after half a decade. Instead, pretraining research grew exponentially in cost, with experiments (not the final training run) dominating the overall expense of model development \citep{cottier2025risingcoststrainingfrontier, epoch2026randdvstrainingcompute}.

Scaling laws failed to deliver because they hit a major roadblock, which is that they are notoriously elusive at small scales. Famously, the two most influential works in the literature reached contradictory conclusions. \citet{Kaplan2020ScalingLaws} held parameters should scale much faster than data, as a result early language models were under-trained \citep{Brown2020Language, rae2022scalinglanguagemodelsmethods, Chowdhery2023PaLM}; \citet{Hoffmann2022AnEmpirical} found parameters and data should scale at the same rate. \citet{Porian2024Resolving} traced this discrepancy to methodological choices such as whether to count embeddings in the parameters, how long to warm up the learning rate, and scale-dependent hyperparameter tuning; choices that disproportionately affect small-scale models \citep{pearce2024reconciling}. This early discord foreshadowed later difficulty. Surveying the literature, \citet{li2025misfitting} found many scaling laws hard to reproduce and, in particular, that scaling laws based on smaller models with 100M parameters or less are unreliable. \citet{Li2025PredictableScale} aptly summarize the situation as the ``\emph{scaling gap} where insights from small-scale experiments often fail to transfer to resource-intensive production systems''.

How do we close this gap and deliver on scaling laws' original promise? We find that the recipe for a scaling law has one key ingredient that, more than any other, determines final quality: rigorous hyperparameter tuning (\S\ref{sec:estimating-small-scale-scaling-laws}). By ablating different choices, we show scaling laws are fairly robust to details like how to count parameters or whether to decay the learning rate (\S\ref{sec:estimating-small-scale-scaling-laws:the-impact-of-methodology}). However, poorly tuned hyperparameters can be disastrous (\S\ref{sec:estimating-small-scale-scaling-laws:the-importance-of-hyperparameter-tuning}). Whereas common practice tunes over a small grid, you may need to search hundreds of configurations to reveal the scaling law in small models. Such a search sounds expensive, yet we also find scaling laws extend remarkably far down; as little as 4M parameters where models train in under an hour on a single GPU, bringing the cost well within reach for many academic and industry labs.

Hyperparameters unduly affect small-scale laws because small-scale models are unduly difficult to tune. We explain this difficulty in terms of the hyperparameter loss surface (\S\ref{sec:understanding-the-hyperparameter-loss-surface-across-scales}). As parameters and data grow together, they drive good configurations to fill a larger fraction of the space and the intrinsic dimension of the surface goes down (\S\ref{sec:understanding-the-hyperparameter-loss-surface-across-scales:hyperparameter-tuning-gets-easier-with-scale}---\S\ref{sec:understanding-the-hyperparameter-loss-surface-across-scales:the-hyperparameter-loss-surface-becomes-lower-dimensional}). Moreveover, as model size increases, the \emph{effective} number of hyperparameters drops to one. This phenomenon introduces a new way of thinking about scaling, where one should do a thorough job sweeping hyperparameters in the smaller scales to due to their sensitivity whereas larger models are easier to tune (\S\ref{sec:understanding-the-hyperparameter-loss-surface-across-scales:implications-for-small-scale-experiments}). Luckily smaller scale experiments are significantly cheaper to run than their larger scale counterparts so proper scaling becomes a matter of proper computational budget allocation. 

Building on our insights, we propose a new methodology for proper scaling drawing on three tools from the literature (\S\ref{sec:designing-small-scale-experiments}). Namely, the \emph{noisy quadratic limit} reveals if hyperparameters are fully tuned and how hard that is to accomplish \citep{lourie2025hyperparameter}; \emph{scaling laws} describe how loss changes as models scale up \citep{rosenfeld2020a, Kaplan2020ScalingLaws, Hoffmann2022AnEmpirical}; and the \emph{perplexity--capability correspondence} answers whether conclusions about pretraining loss have consequences for the capabilities we ultimately care about \citep{mayilvahanan25a}.\footnote{
    Even when we often can not predict these capabilities in advance \citep{lourie-etal-2025-scaling}.
} Small-scale experiments exploit strong assumptions and we argue for an approach where each assumption becomes a diagnostic that improves our qualitative understanding of how the model scales (\ref{sec:designing-small-scale-experiments:a-methodology-for-small-scale-experiments}). This qualitative understanding is important because even though phenomena like scaling laws exist at the small scale, quantitative techniques such as extrapolating the loss compound small errors (\ref{sec:designing-small-scale-experiments:the-pitfalls-of-pure-extrapolation}). Finally, we apply our takeaways to a question that once took the field years to settle (\ref{sec:designing-small-scale-experiments:case-study-pre-norm-vs-post-norm-architectures}): where normalization layers should go in the transformer architecture, and find that larger scales favor pre-normalization \citep{Vaswani2017Attention, baevski2018adaptiveinputrepresentationsneural, Xiong2020OnLayer}.


In the end, these ideas give a practical methodology for model-centric research, which  improves the model holding \textit{data fixed} in contrast to \emph{data-centric} research, which improves the data holding the model fixed. 
The methodology is as follows: first build a complete understanding at the small scale by thoroughly exploring the model's hyperparameters, then use simple rules to carry it up. This regimen offers the rare opportunity to cut costs and raise rigor, but it takes the right tools discussed in this work. With them, small-scale experiments can finally deliver on their long-awaited promise.
We note finally that our discussion focuses on model-centric research while data-centric research likely requires different techniques (\S\ref{sec:conclusion}). 